% Standalone appendix fragment; amsmath, amssymb, amsthm required.
\section{What excess Gaussian energy does and does not control}
\label{app:excessenergy}

Write a width-$N$, depth-$L$ tanh network as
$h^1=\tanh(W^1\bar x)$, $h^s=\tanh(W^sh^{s-1})$ for $2\le s<L$,
and $z=v^Th^{L-1}$, where $\bar x=(1,x)$ and $x\in\{-1,1\}^d$.
Its weighted Gaussian energy is
\[
 \kappa_N=\frac dN\|W^1\|_F^2+
       \sum_{s=2}^{L-1}\|W^s\|_F^2+N\|v\|_2^2.
\]
The following facts concern this actual energy. They do not establish a
sample-complexity lower bound for the unrestricted Gaussian posterior.

\begin{proposition}[Excess energy controls total deviation from rank one]
Suppose $|z(x)|\ge1/2$ on the whole Boolean cube. Let $R_s$ be a best
rank-one Frobenius approximation to $W^s$, and let
$\rho_s=\|W^s\|_F^2/\|W^s\|_{\mathrm{op}}^2$ be its stable rank.
Then
\begin{align}
 \sum_{s=2}^{L-1}\|W^s-R_s\|_F^2
 &\le \kappa_N-L(d/4)^{1/L},\label{eq:excessrank}\\
 \sum_{s=2}^{L-1}\log\rho_s
 &\le L\log(\kappa_N/L)-\log(d/4).
\end{align}
In particular, if $\kappa_N\le L+B$, both right-hand sides are at most
$B-\log(d/4)$.
\end{proposition}

\begin{proof}
Set $a=\|W^1\|_F/\sqrt N$, $c=\sqrt N\|v\|_2$, and
$p_s=\|W^s\|_{\mathrm{op}}$. Cube averaging supplies an input with
$\|W^1\bar x\|_2/\sqrt N\le a$. Tanh contraction gives
$1/2\le ac\prod_s p_s$. Thus the product of the $L$ numbers
$da^2,p_2^2,\ldots,p_{L-1}^2,c^2$ is at least $d/4$, and their sum is at
least $L(d/4)^{1/L}$. Subtract this sum from $\kappa_N$ and use
$\|W^s-R_s\|_F^2=\|W^s\|_F^2-p_s^2$ to obtain
\eqref{eq:excessrank}. Alternatively, the product of the Frobenius energy
terms is at least $(d/4)\prod_s\rho_s$ and at most $(\kappa_N/L)^L$,
which proves the logarithmic bound. Finally use $e^u\ge1+u$ and
$\log(1+u)\le u$.
\end{proof}

Thus bounded excess permits only a bounded total squared residual after
replacing internal matrices by rank-one matrices. This does not justify
replacing those matrices in the computation: the suffix can amplify the
replacement errors. The next example shows that even a constant excess
allowance does not imply inverse-polynomial hidden-state robustness.

\begin{proposition}[Bounded excess permits a fragile scalar transmission]
For $d=4$, every even $L\ge4$, and every $N\ge1$, there is a network with
$z(x)=1$ for all inputs and
\[
 \kappa_N\le L+76\le K_*(4,L)+76,
\]
whose output can be made zero by a perturbation at one hidden layer of
normalized RMS at most $L^{-1}e^{-\sqrt L/4}$.
Here $K_*(4,L)$ is the scalar minimum for margin-$1/2$ transmission
defined above. It also equals the globally minimum margin-$1/2$ energy
for this model, by the simple-target result in the research report.
\end{proposition}

\begin{proof}
Put $n=L/2$, $h=L^{-1/2}$, $k_s=\min(s,L-s)$, and
\[
 r_s=\frac{e^{-h(k_s-1)}}{2k_s},\qquad 1\le s<L.
\]
Give each first-layer row the sole constant-input weight
$\operatorname{atanh}(1/2)$. Let every entry of $W^s$ be $b_s/N$, where
$b_s=\operatorname{atanh}(r_s)/r_{s-1}$, and let $v_i=2/N$.
The hidden vector is exactly $r_s\mathbf1$, so the output is
$2r_{L-1}=1$ and the energy is
\[
 4\operatorname{atanh}(1/2)^2+\sum_{s=2}^{L-1}b_s^2+4.
\]
For completeness, the following loose constants are uniform in depth.
The ratios $r_s/r_{s-1}$ occur in reciprocal pairs $t_k,t_k^{-1}$,
where $t_k=(1+1/k)e^h$ and $1\le k<n$. Since $h\le1/2$,
\[
 t_k-t_k^{-1}\le e^h(2/k+2h),\qquad
 t_k^2+t_k^{-2}-2\le8e^{2h}(k^{-2}+h^2).
\]
Using $e^{2h}<3$, $\sum_{k\ge1}k^{-2}<2$, and $(n-1)h^2\le1/2$ gives
$\sum_s(r_s/r_{s-1})^2<L-2+60$.
Also $\max_s(r_s/r_{s-1})^2<12$ and $\sum_s r_s^2<1$.
For $0<t\le1/2$ the power series gives
$(\operatorname{atanh}(t)/t)^2\le1+t^2$. Consequently
\[
 \sum_s b_s^2<L-2+72,
 \qquad 4\operatorname{atanh}(1/2)^2\le5/4,
 \qquad \kappa_N<L+76.
\]
The baseline lower bound gives $K_*(4,L)\ge L$.
Finally, adding $-r_n\mathbf1$ after the middle activation zeros that
state and hence every later state. Its normalized RMS is
\[
 r_n=L^{-1}\exp\!\left(-\frac{L/2-1}{\sqrt L}\right)
 \le L^{-1}e^{-\sqrt L/4}.
\]
\end{proof}

This example represents a statistically easy target: the report's
simple-target theorem gives depth-independent strong-prior learning
for constants, at fixed width and depth before taking the variance limit. It therefore
separates fragility of a cheap parameter realization from difficulty of
learning its Boolean function. It does not show that the fragile realization
is a minimum-energy representative.

\begin{proposition}[Additional depth has asymptotically unit energy cost]
Suppose a width-$N$, depth-$L_0$ network realizes $f$ with margin at least
$1/2$. Write $c^2=N\|v\|_2^2>0$ and $E_{\mathrm{pre}}=\kappa_N-c^2$.
For every integer $t\ge1$, there is a width-$N$, depth-$(L_0+t)$ realization
of the same function, with margin at least $1/2$ and energy at most
\[
 E_{\mathrm{pre}}+(t+13)\max\{c^2,1\}^{1/(t+1)}.
\]
Write $\kappa_{N,L}(f)$ for the infimum energy of depth-$L$, width-$N$
representations with signed output strictly greater than $1/2$ on the cube.
For each fixed representable $f,d,N$, this infimum satisfies
$\kappa_{N,L}(f)=L+O_{f,d,N}(1)$ as $L\to\infty$.
\end{proposition}

\begin{proof}
Extend the scalar minimum $K_*(D,m)$ to any real $D>0$: minimize
$Da^2+\sum_{j=2}^{m-1}b_j^2+c'^2$ subject to $c'r_{m-1}\ge1/2$,
where $r_1=\tanh a$ and $r_j=\tanh(b_jr_{j-1})$.
The scalar baseline upper-bound argument works unchanged for real $D>0$:
\[
 K_*(D,m)\le(m+12)(\max\{D,4\}/4)^{1/m}.
\]
Indeed, setting $q=(\max\{D,4\}/4)^{1/(2m)}$ and
$r_j=q^j/[\sqrt{\max\{D,4\}}\min(j,m-j)]$ gives this bound by
$a=\operatorname{atanh}r_1$, $b_j=\operatorname{atanh}(r_j)/r_{j-1}$,
$c'=q$, and the summable reciprocal-ratio estimates used earlier.

Take the scalar tuple for $D=4c^2$ and $m=t+1$. Replace the old readout
by a hidden layer whose every row is $2a v^T$, followed by the remaining
scalar layers and readout. Its first new preactivation is $2az(x)$, whose
absolute value is at least $a$ and whose sign is $f(x)$. Monotonicity and
oddness give the same sign with final margin at least $1/2$.
The first new block has energy $4a^2c^2$; the remaining new blocks have
exactly the remaining scalar costs. The total is therefore
$E_{\mathrm{pre}}+K_*(4c^2,t+1)$, proving the displayed upper bound.
An arbitrarily small readout increase supplies strict margins.
Expanding this bound as $t\to\infty$, and combining it with
$L(d/4)^{1/L}\le\kappa_{N,L}(f)$, proves the final assertion.
\end{proof}

These energy estimates alone do not settle unrestricted Gaussian sample
complexity. They rule out exponential growth of minimum target energy under depth
padding, while showing that bounded excess does not make every cheap
realization robust. A target-relative argument must also control competing
functions and their prior masses; subtracting a common transmission cost
alone does not do so.
